\section{本章小结}
\label{sec:lfunc-summary}

本章开头问：哪些标量读数能由函数的范数稳定控制，它们长什么样，能否从局部扩展到全空间。泛函是从线性空间到标量的映射；加权积分在等价类上有定义，点评估只在由真正的函数组成的空间上有定义，平方范数则是非线性的。线性泛函连续当且仅当存在常数$C$使$\lvert\ell(x)\rvert\leq C\lVert x\rVert$，最小的$C$就是泛函范数。同一条读数规则是否连续、范数多大，取决于所配的范数：点评估在一致范数下范数为$1$，在积分范数下被窄尖峰破坏而不连续。

全部连续线性泛函组成对偶空间，无论原空间是否完备，对偶空间都是Banach空间，它的对偶是二次对偶。Hilbert空间中，每个连续线性泛函都是与唯一向量的内积，范数相等，证明只用正交投影而不要求可分；表示的前提是泛函已经连续、空间已经完备。$\sigma$有限测度下$1\leq p<\infty$时，$L^p$上的连续线性泛函都是以$L^q$函数为权重的积分，权重在几乎处处意义下唯一、范数相等；$p=\infty$时对偶更大。

一般空间中写不出表示，Hahn--Banach定理转而保证子空间上的有界线性泛函能保持范数地延拓到全空间，延拓一般不唯一。由此每个向量的范数等于单位泛函读数的最大值，不同向量能被连续线性读数区分，闭子空间外的点能被一个在子空间上为零的读数分开；$\ell^\infty$上还存在不来自$\ell^1$的连续线性泛函。完备性给出另一种统一控制：Banach空间上逐点有界的泛函族范数一致有界，不完备时可以失败。

弱收敛要求每个连续线性读数收敛，弱星收敛只要求在原空间的每个向量上求值收敛，二者的观测来源不同。范数收敛推出弱收敛，逆向在无限维中不成立，标准正交序列与快速振荡都弱收敛到零而范数不变；只检查一部分读数不能代替全部读数。对偶空间中的弱收敛推出弱星收敛，逆向需要二次对偶中没有多余元素。弱收敛列有界、范数下半连续，Banach前对偶上的弱星收敛列也有界；有界只是必要条件。

紧性方面，任意赋范空间的对偶闭球在弱星拓扑下紧；反射空间通过等距同构$J$把这一结论转回原空间，闭球弱紧。紧性要转成子列，还需要条件：前对偶可分时，对偶闭球上的弱星拓扑可度量，有界序列有弱星收敛子列；反射空间中由Eberlein--\v{S}mulian定理，有界序列有弱收敛子列，不需可分性。$L^1$中质量集中的序列说明非反射空间里这可以失败。这些都不是范数紧性。

本章的读数都输出一个标量。第~\ref{chap:linear-operators}章把输出推广为向量，研究有界线性算子、算子范数、伴随与完备空间中的算子定理；本章的对偶空间、Riesz表示和弱收敛，是那里定义伴随与讨论紧算子的工具。
